\subsection{集合可积性的判别法}

命题 10. --- 为使 X 中的开（相应地，闭）集 A 可积，必须且只须 \(|\mu|^{*}(A)<+\infty\) 。

由于此时 \(\varphi_{A}\) 是下半（相应地，上半）连续的，命题由 No.~4 的命题 5 及其推论 1 得出。

推论 1. --- 每个紧集都是可积的；每个相对紧的开集都是可积的。

推论 2. --- 对 X 上的每个正测度 \(\mu\) ，A \(\mapsto\) \(\mu^{*}(\mathrm{A})\) 是 X 上的一个容量（参见 GT, IX, §6, No.~9, 例子）。

例子. --- 对 R 上的 Lebesgue 测度 \(\mu\) ，由命题 10 知每个有界开区间 {]}a,b{[} 都是可积的且测度为 b-a（§1, No.~2, 命题 9）。由于每个退缩为一点的集对 Lebesgue 测度都是可忽略的，于是以 a 与 b 为端点的所有区间都有相同的测度 b-a。

命题 11. --- 设 \(\mathfrak{G}\) 是 X 中可积开集所成的一个集，按关系 \(\subset\) 有向；为使 \(A = \bigcup_{G \in \mathfrak{G}} G\) 可积，必须且只须 \(\sup_{G \in \mathfrak{G}} |\mu|(G) < +\infty\) ；此时 \(\mu(A) = \lim_{\mathfrak{G}} \mu(G)\) 且 \(|\mu|(A) = \sup_{G \in \mathfrak{G}} |\mu|(G)\) 。

因为，已知 \(|\mu|^*(\mathbf{A}) = \sup_{\mathbf{G} \in \mathfrak{G}} |\mu|(\mathbf{G})\) （§1, No.~2, 命题 7）；于是命题由命题 10 得出。

推论. --- 设 \(\mathfrak{F}\) 是 X 中可积闭集所成的一个集，按关系 \(\supset\) 有向；则闭集 \(B = \bigcap_{H \in \mathfrak{F}} H\) 是可积的，且有 \(\mu(B) = \lim_{\mathfrak{F}} \mu(H)\) 与 \(|\mu|(B) = \inf_{H \in \mathfrak{F}} |\mu|(H)\) 。

因为，设 \(H_{0}\) 是 F 中的一个集；由于 \(H_{0}\) 可积，它含于一个可积开集 U 中（§1, No.~4, 命题 19）；开集 \(U \cap CH\) 构成一个按关系 \(\subset\) 有向的集，它们都含于 U 中，且并集为 \(U \cap CB\) ；于是我们归结为命题 11。

定理 4. --- 为使一个集 A 可积，必须且只须：对每个 \(\varepsilon > 0\) ，存在可积开集 G 与紧集 K，使 \(K \subset A \subset G\) 且

\[
| \mu | (\mathrm{G} - \mathrm{K}) = | \mu | (\mathrm{G}) - | \mu | (\mathrm{K}) \leqslant \varepsilon .
\]

\begin{enumerate}
\def\labelenumi{\alph{enumi})}
\item
  条件是充分的，因为它意味着 \(\varphi_{\mathrm{K}} \leqslant \varphi_{\mathrm{A}} \leqslant \varphi_{\mathrm{G}}\) 且 \(\int (\varphi_{\mathrm{G}} - \varphi_{\mathrm{K}}) d|\mu| \leqslant \varepsilon\) ；由于 \(\varphi_{\mathrm{G}}\) 与 \(\varphi_{\mathrm{K}}\) 可积，\(\varphi_{\mathrm{A}}\) 也可积（§3, No.~4, 命题 8）。
\item
  条件是必要的。若 A 可积，则存在开集 \(G \supset A\) 使 \(|\mu|^{*}(G)\) 可任意接近 \(|\mu|^{*}(A) = |\mu|(A)\) （§1, No.~4, 命题 19）；于是，一切归结为证明：对每个 \(\varepsilon > 0\) ，存在紧集 \(K \subset A\) 使 \(|\mu|(A) - |\mu|(K) \leqslant \varepsilon\) 。由于 \(\varphi_{A}\) 可积，存在函数 \(f \geqslant 0\) ，上半连续且具有紧支集 S，使 \(f \leqslant \varphi_{A}\) 且 \(\int(\varphi_{A} - f)d|\mu| \leqslant \varepsilon/2\) （No.~4, 定理 3）。设 \(\delta > 0\) 是任意一个数，并设 K 是点 \(x \in X\) （使 \(f(x) \geqslant \delta\) 的）所成之集；K 是闭的且含于 S 中，因而是紧的，且由于 \(f \leqslant \varphi_{A}\) ，我们有 \(K \subset A\) 。集 B = A - K 是可积的，且 \(f \leqslant \varphi_{K} + \delta\varphi_{B}\) ，由此
\end{enumerate}

\[
\int f d | \mu | \leqslant | \mu | (\mathrm{K}) + \delta \cdot | \mu | (\mathrm{B}) \leqslant | \mu | (\mathrm{K}) + \delta \cdot | \mu | (\mathrm{A}),
\]

最后

\[
| \mu | (\mathrm{A}) \leqslant \int f d | \mu | + \frac {\varepsilon}{2} \leqslant | \mu | (\mathrm{K}) + \delta \cdot | \mu | (\mathrm{A}) + \frac {\varepsilon}{2},
\]

这就完成了证明，因为 \(\delta\) 是任意的。

推论 1. --- 为使一个集 A 可积，必须且只须：对每个 \(\varepsilon > 0\) ，存在紧集 \(K \subset A\) 使 \(|\mu|^*(A - K) \leqslant \varepsilon\) 。此时测度 \(|\mu|(A)\) 是诸测度 \(|\mu|(K)\) （紧集 \(K \subset A\) 的）所成之集的上确界。

条件是必要的，因为若 G 与 K 满足定理 4 的条件，则 \(|\mu|^{*}(A - K) \leqslant |\mu|^{*}(G - K) \leqslant \varepsilon\) 。

条件是充分的，因为它表明：对于平均收敛拓扑，\(\varphi_{A}\) 处于可积函数 \(\varphi_{K}\) （K 为 A 的任意紧子集）所成之集的闭包之中。

推论 2. --- 对每个可积集 A，存在：

\(1^{\circ}\) 一个集 \(\mathbf{A}_1\supset \mathbf{A}\) ，它是可积开集的可数交，使 \(\mathbf{A}_1 - \mathbf{A}\) 是可忽略的；

\(2^{\circ}\) 一个集 \(\mathbf{A}_2\subset \mathbf{A}\) ，它是两两不相交的紧集的可数并，使 \(\mathbf{A} - \mathbf{A}_2\) 是可忽略的。

\(1^{\circ}\) 对每个整数 \(n\) ，存在可积开集 \(\mathbf{G}_n \supset \mathbf{A}\) 使 \(|\mu| (\mathrm{G}_n) - |\mu| (\mathrm{A}) \leqslant 1/n\) ；若 \(\mathbf{A}_1\) 是诸 \(\mathbf{G}_n\) 的交，则 \(|\mu| (\mathrm{A}_1) = |\mu| (\mathrm{A})\) （No.~5, 命题 7 的推论），因此 \(\mathbf{A}_1 - \mathbf{A}\) 是可忽略的。

\(2^{\circ}\) 我们按下述方式归纳地定义紧集 \(\mathbf{K}_n\) ：\(\mathbf{K}_1 \subset \mathbf{A}\) 且 \(|\mu| (\mathrm{A} - \mathrm{K}_1) \leqslant 1\) ；\(\mathbf{K}_n \subset \mathbf{A} - \bigcup_{i=1}^{n-1} \mathbf{K}_i\) 且 \(|\mu| (\mathrm{A} \cap \mathbb{C}(\bigcup_{i=1}^{n-1} \mathbf{K}_i) \cap \mathbb{C} \mathbf{K}_n) \leqslant 1/n\) 对 \(n > 1\) （定理 4）；若 \(\mathbf{A}_2\) 是诸 \(\mathbf{K}_n\) 的并，则 \(|\mu| (\mathrm{A}_2) = |\mu| (\mathrm{A})\) （No.~5, 命题 8），因此 \(\mathbf{A} - \mathbf{A}_2\) 是可忽略的。

推论 3. --- 每个外测度有限的集都含于一个可忽略集与一个两两不相交紧集的可数族（其测度之和有限）的并中。

只需把推论 2 应用于包含所给集的一个可积开集。

推论 4. --- 对 X 中每个开集 U，\(|\mu|^{*}(U)\) 是诸测度 \(|\mu|(K)\) （紧集 \(K \subset U\) 的）的上确界。

若 \(|\mu|^*(\mathrm{U}) < +\infty\) ，这由定理 4 立得。下述论证也涵盖 \(|\mu|^*(\mathrm{U}) = +\infty\) 的情形。由于 X 是局部紧的且 U 是开的，\(\varphi_{\mathrm{U}}\) 是函数 \(f \in \mathcal{K}_{+}\) （使 \(f \leqslant \varphi_{\mathrm{U}}\) 且 \(\operatorname{Supp}(f) \subset \mathrm{U}\) 的）所成之集 H 的上包络（参见 §1, No.~1, 引理的证明），且由于 H 按 \(\leqslant\) 有向，我们有 \(|\mu|^*(\mathrm{U}) = \sup_{f \in \mathrm{H}} |\mu|(f)\) （§1, No.~1,

定理 1）；于是由下述事实立得推论：若 \(f \in \mathbf{H}\) 且 \(K = \operatorname{Supp}(f)\) ，则 \(f \leqslant \varphi_{K} \leqslant \varphi_{U}\) 。

注意，\(|\mu|^*(\mathrm{U})\) 也是诸测度 \(|\mu|(\mathrm{G})\) （相对紧开集的，使 \(\overline{\mathrm{G}} \subset \mathrm{U}\) ）的上确界。因为，若 \(\mathrm{K}\) 是含于 \(\mathrm{U}\) 的一个紧集，则对每个 \(x \in \mathrm{K}\) ，存在一个相对紧开邻域 \(\mathrm{V}\) （\(x\) 的）使 \(\overline{\mathrm{V}} \subset \mathrm{U}\) 。用有限个这样的邻域覆盖 \(\mathrm{K}\) ，它们的并 \(\mathrm{G}\) 是一个相对紧开集，使 \(\overline{\mathrm{G}} \subset \mathrm{U}\) 且 \(\mathrm{K} \subset \mathrm{G}\) ，由此 \(|\mu|(\mathrm{K}) \leqslant |\mu|(\mathrm{G}) \leqslant |\mu|^*(\mathrm{U})\) 。

